% Options for packages loaded elsewhere
\PassOptionsToPackage{unicode}{hyperref}
\PassOptionsToPackage{hyphens}{url}
\documentclass[
  11pt,
]{article}
\usepackage{xcolor}
\usepackage[margin=1in]{geometry}
\usepackage{amsmath,amssymb}
\setcounter{secnumdepth}{-\maxdimen} % remove section numbering
\usepackage{iftex}
\ifPDFTeX
  \usepackage[T1]{fontenc}
  \usepackage[utf8]{inputenc}
  \usepackage{textcomp} % provide euro and other symbols
\else % if luatex or xetex
  \usepackage{unicode-math} % this also loads fontspec
  \defaultfontfeatures{Scale=MatchLowercase}
  \defaultfontfeatures[\rmfamily]{Ligatures=TeX,Scale=1}
\fi
\usepackage{lmodern}
\ifPDFTeX\else
  % xetex/luatex font selection
\fi
% Use upquote if available, for straight quotes in verbatim environments
\IfFileExists{upquote.sty}{\usepackage{upquote}}{}
\IfFileExists{microtype.sty}{% use microtype if available
  \usepackage[]{microtype}
  \UseMicrotypeSet[protrusion]{basicmath} % disable protrusion for tt fonts
}{}
\makeatletter
\@ifundefined{KOMAClassName}{% if non-KOMA class
  \IfFileExists{parskip.sty}{%
    \usepackage{parskip}
  }{% else
    \setlength{\parindent}{0pt}
    \setlength{\parskip}{6pt plus 2pt minus 1pt}}
}{% if KOMA class
  \KOMAoptions{parskip=half}}
\makeatother
\setlength{\emergencystretch}{3em} % prevent overfull lines
\providecommand{\tightlist}{%
  \setlength{\itemsep}{0pt}\setlength{\parskip}{0pt}}
\usepackage{bookmark}
\IfFileExists{xurl.sty}{\usepackage{xurl}}{} % add URL line breaks if available
\urlstyle{same}
\hypersetup{
  pdftitle={Writing AI-Assisted Textbooks},
  hidelinks,
  pdfcreator={LaTeX via pandoc}}

\title{Writing AI-Assisted Textbooks}
\author{}
\date{}

\begin{document}
\maketitle

\section{is there a paper about how to write an AI assisted textbooks}\label{is-there-a-paper-about-how-to-write-an-ai-assisted-textbooks}

There isn't (yet) a definitive ``how to write an academic textbook with AI'' manual, but there is emerging practice from textbook authors and AI users that's mature enough to turn into a workable workflow.\footnote{https://www.youtube.com/watch?v=DFeP863BaPM}\footnote{https://community.openai.com/t/technique-for-writing-entire-books/705519}\footnote{https://blog.taaonline.net/2023/11/textbook-and-academic-authors-share-how-they-use-generative-ai/}\footnote{https://www.reddit.com/r/OpenAI/comments/1cgxtdd/using\_ai\_to\_write\_a\_book/}

Below is a concrete, academically-defensible way to use AI as a coauthoring tool rather than a ghostwriter.

\subsection{Overall workflow}\label{overall-workflow}

\begin{enumerate}
\def\labelenumi{\arabic{enumi}.}
\tightlist
\item
  Define the book and constraints yourself

  \begin{itemize}
  \tightlist
  \item
    Specify course, level, pre-reqs, learning objectives, and your distinctive angle; this is human work, not AI.
  \item
    Treat AI as an assistant that operates inside those constraints (like a smart RA), not as the source of the intellectual contribution.\footnote{https://community.openai.com/t/technique-for-writing-entire-books/705519}\footnote{https://blog.taaonline.net/2023/11/textbook-and-academic-authors-share-how-they-use-generative-ai/}
  \end{itemize}
\item
  Work chapter-by-chapter in small chunks

  \begin{itemize}
  \tightlist
  \item
    A common pattern is: you design the outline, then have AI help generate or refine sections, examples, and exercises one at a time.\footnote{https://www.reddit.com/r/OpenAI/comments/1cgxtdd/using\_ai\_to\_write\_a\_book/}\footnote{https://community.openai.com/t/technique-for-writing-entire-books/705519}
  \item
    Avoid ``write the whole textbook'' prompts; they give shallow, repetitive prose and raise originality/plagiarism concerns.\footnote{https://community.openai.com/t/technique-for-writing-entire-books/705519}\footnote{https://www.reddit.com/r/OpenAI/comments/1cgxtdd/using\_ai\_to\_write\_a\_book/}
  \end{itemize}
\item
  Keep a human‑curated backbone

  \begin{itemize}
  \tightlist
  \item
    Maintain a master outline, style guide, and list of key references that you update as you go; feed these repeatedly into the model so it stays on-pitch.\footnote{https://www.reddit.com/r/OpenAI/comments/1cgxtdd/using\_ai\_to\_write\_a\_book/}\footnote{https://community.openai.com/t/technique-for-writing-entire-books/705519}
  \item
    Think of this as your ``corpus'' that the model is conditioned on, not something the model invents.
  \end{itemize}
\end{enumerate}

\subsection{Where AI is most useful for textbook writing}\label{where-ai-is-most-useful-for-textbook-writing}

Based on what textbook authors report actually using generative models for, plus what general ``AI book writing'' workflows emphasize, the sweet spots are:\footnote{https://blog.taaonline.net/2023/11/textbook-and-academic-authors-share-how-they-use-generative-ai/}\footnote{https://community.openai.com/t/technique-for-writing-entire-books/705519}\footnote{https://www.reddit.com/r/OpenAI/comments/1cgxtdd/using\_ai\_to\_write\_a\_book/}

\begin{itemize}
\tightlist
\item
  Brainstorming and refining structure

  \begin{itemize}
  \tightlist
  \item
    Ask for alternative chapter organizations, prerequisite maps, and concept progressions once you've given it your first-pass outline.\footnote{https://community.openai.com/t/technique-for-writing-entire-books/705519}\footnote{https://www.reddit.com/r/OpenAI/comments/1cgxtdd/using\_ai\_to\_write\_a\_book/}
  \item
    Use it to identify missing topics or perspectives in a draft TOC: ``Given this outline for an insurance regulation text, what major regulatory issues might be missing?''\footnote{https://blog.taaonline.net/2023/11/textbook-and-academic-authors-share-how-they-use-generative-ai/}
  \end{itemize}
\item
  Drafting neutral, expository prose from your notes

  \begin{itemize}
  \tightlist
  \item
    Many authors dictate or rough-draft a section, then feed transcripts/notes to AI to turn into a clean first draft, which they then edit heavily.\footnote{https://www.youtube.com/watch?v=DFeP863BaPM}\footnote{https://www.reddit.com/r/OpenAI/comments/1cgxtdd/using\_ai\_to\_write\_a\_book/}
  \item
    You can keep control of all definitions, notation, and key arguments but offload smoothing, clarification, and example elaboration.
  \end{itemize}
\item
  Generating pedagogical scaffolding

  \begin{itemize}
  \tightlist
  \item
    Authors report using AI to propose learning objectives, key terms, chapter summaries, and discussion questions based on an existing chapter draft.\footnote{https://blog.taaonline.net/2023/11/textbook-and-academic-authors-share-how-they-use-generative-ai/}
  \item
    You then vet these for correctness, level, and alignment with your course.
  \end{itemize}
\item
  Language polishing and consistency

  \begin{itemize}
  \tightlist
  \item
    Many academics use AI for grammar, simplification, and style uniformity across chapters, especially in multi-author volumes.\footnote{https://blog.taaonline.net/2023/11/textbook-and-academic-authors-share-how-they-use-generative-ai/}
  \item
    You can also use it to enforce a consistent voice or complexity level by giving it your style guide and one or two ``gold standard'' passages.
  \end{itemize}
\item
  Marketing and ancillaries

  \begin{itemize}
  \tightlist
  \item
    Authors are using AI to help with book proposals, blurbs, slide outlines, and even marketing copy for the text.\footnote{https://blog.taaonline.net/2023/11/textbook-and-academic-authors-share-how-they-use-generative-ai/}
  \end{itemize}
\end{itemize}

\subsection{A concrete AI‑assisted textbook workflow}\label{a-concrete-aiassisted-textbook-workflow}

Here is a practical pipeline you could use for, say, a graduate insurance regulation text.

\begin{enumerate}
\def\labelenumi{\arabic{enumi}.}
\tightlist
\item
  Design the intellectual skeleton (human)

  \begin{itemize}
  \tightlist
  \item
    Write a 1--2 page concept note: target readers, core thesis of the book, course use cases, and what's distinctive about your approach.
  \item
    Draft a detailed TOC: parts, chapters, and 3--7 bullet points per chapter.
  \end{itemize}
\item
  Use AI to stress‑test and refine the outline

  \begin{itemize}
  \tightlist
  \item
    Prompt: ``Here is my outline for a graduate textbook on insurance regulation in the US and EU. Identify gaps, redundancies, and alternative structures, but do not invent new doctrinal claims.''
  \item
    Integrate only those suggestions that fit your own understanding; do not adopt legal or regulatory claims without checking them.
  \end{itemize}
\item
  Create a living style and pedagogy guide

  \begin{itemize}
  \tightlist
  \item
    Decide on: notation, citation style, voice (first/third person), depth of math, typical chapter structure (e.g., intro, conceptual exposition, case, exercises).
  \item
    Feed this plus one fully human‑written sample section (e.g., your best lecture notes turned into prose) and tell the model: ``This is the target style. When drafting or editing, match this.''
  \end{itemize}
\item
  Drafting chapters using ``notes → AI draft → human'' loops For each major section:

  \begin{itemize}
  \tightlist
  \item
    You write structured notes, bullet lists, and references (including specific statutes, cases, or empirical papers).
  \item
    AI step: ``Turn these notes into a clear, 1500-word expository section for graduate students, preserving all technical details and citations, avoiding invented cases or laws.''
  \item
    You then:

    \begin{itemize}
    \tightlist
    \item
      Check every legal, regulatory, and empirical claim against primary sources.
    \item
      Rewrite where nuance or originality is crucial (e.g., your own empirical findings or doctrinal argument).
    \end{itemize}
  \end{itemize}
\item
  Generating exercises and ancillaries

  \begin{itemize}
  \tightlist
  \item
    Use AI to propose end-of-chapter questions, short problems, and ``what if'' policy scenarios based on your draft.\footnote{https://blog.taaonline.net/2023/11/textbook-and-academic-authors-share-how-they-use-generative-ai/}
  \item
    You filter for correctness and difficulty, and make sure any data or statutes are real and properly dated.
  \item
    For slides and lecture outlines, you can feed a chapter and ask for a 20-slide outline with key diagrams and examples; then you reconstruct the actual slides yourself.
  \end{itemize}
\item
  Managing long context and coherence

  \begin{itemize}
  \tightlist
  \item
    Some authors build tools where the AI always sees a JSON-like tree of Book → Chapter → Section, and content is generated per node, not all at once.\footnote{https://community.openai.com/t/technique-for-writing-entire-books/705519}
  \item
    Even without custom tools, you can do something similar by:

    \begin{itemize}
    \tightlist
    \item
      Keeping per-chapter summaries that you paste into new prompts.
    \item
      Periodically asking the model: ``Given chapter summaries 1--4, draft a 2-page overview that shows how they fit together'' to detect coherence issues.
    \end{itemize}
  \end{itemize}
\end{enumerate}

\subsection{Guardrails: academic integrity and quality}\label{guardrails-academic-integrity-and-quality}

\begin{itemize}
\tightlist
\item
  No AI‑only doctrinal statements or citations

  \begin{itemize}
  \tightlist
  \item
    Treat any legal citation, numeric claim, or literature reference from AI as a lead to be verified, never as authoritative.
  \item
    Maintain your own BibTeX / reference list built from actual databases; use AI at most to format or group them.
  \end{itemize}
\item
  Transparency and authorship

  \begin{itemize}
  \tightlist
  \item
    Many textbook authors now disclose in their prefaces that they used generative tools for drafting, editing, or ancillary materials, while emphasizing that all intellectual content and responsibility remain with them.\footnote{https://blog.taaonline.net/2023/11/textbook-and-academic-authors-share-how-they-use-generative-ai/}
  \item
    This also mitigates concerns from publishers and adopters.
  \end{itemize}
\item
  Copyright and originality

  \begin{itemize}
  \tightlist
  \item
    Avoid asking the model to mimic particular proprietary textbooks or to ``rewrite chapter 3 of X.''
  \item
    Work from your own notes, slides, and articles; AI then helps reorganize and phrase, which is far less likely to generate infringing or derivative text.
  \end{itemize}
\item
  Data and privacy

  \begin{itemize}
  \tightlist
  \item
    If you include proprietary data or in-progress empirical results, use a deployment that your institution approves for research data (or keep those sections human-written and only use AI for surrounding explanation).
  \end{itemize}
\end{itemize}

\subsection{Tools and setups people actually use}\label{tools-and-setups-people-actually-use}

\begin{itemize}
\tightlist
\item
  General-purpose models (like this one) for brainstorming, outlining, editing, and question generation --- usually via chat interfaces.\footnote{https://www.reddit.com/r/OpenAI/comments/1cgxtdd/using\_ai\_to\_write\_a\_book/}\footnote{https://blog.taaonline.net/2023/11/textbook-and-academic-authors-share-how-they-use-generative-ai/}
\item
  Note‑to‑draft workflows where authors record themselves talking through a chapter, transcribe, and then have AI convert the transcript to prose.\footnote{https://www.reddit.com/r/OpenAI/comments/1cgxtdd/using\_ai\_to\_write\_a\_book/}
\item
  Author‑built apps that represent the book as a tree and generate content per node, often using retrieval-augmented generation to keep previous chapters in view.\footnote{https://community.openai.com/t/technique-for-writing-entire-books/705519}
\item
  Off-the-shelf ``AI book writer'' tools exist, but most are tuned for trade books and speed, not scholarly rigor; they're best avoided or used only for very preliminary brainstorming.\footnote{https://www.inkfluenceai.com/ai-book-writer}\footnote{https://www.youtube.com/watch?v=DFeP863BaPM}
\end{itemize}

If you tell me the specific course you have in mind (e.g., grad insurance regulation vs ERM for MBAs), I can sketch an AI‑integrated project plan tailored to that book, including sample prompts for one concrete chapter.

\end{document}
